%% Appendix D: Detailed Generalization Evidence
%% Complete methods and results for the 13 cross-participant, 2 cross-session, 2 cross-language, and 5 cross-modal studies

\subsection{Cross-Participant Generalization Studies (n=13)}

[PLACEHOLDER: For each of the 13 studies demonstrating cross-participant generalization:]
\begin{itemize}
    \item Study title and authors
    \item Participant count (training vs. test)
    \item Transfer learning methodology
    \item Performance on training participant(s)
    \item Performance on novel participant(s)
    \item Performance retention (\%)
    \item Neural features or mechanisms underlying transfer
\end{itemize}

\subsection{Cross-Session Generalization Studies (n=2)}

[PLACEHOLDER: Detailed analysis of the two studies evaluating cross-session stability]
\begin{itemize}
    \item Time span between sessions
    \item Performance retention over time
    \item Apparent sources of drift (electrode drift, neural plasticity, etc.)
    \item Adaptation/recalibration requirements
\end{itemize}

\subsection{Cross-Language Generalization Studies (n=2)}

[PLACEHOLDER: Analysis of cross-language transfer]
\begin{itemize}
    \item Languages tested
    \item Training language(s) and test language(s)
    \item Performance on trained language
    \item Performance on novel language
    \item Linguistic properties facilitating or hindering transfer
\end{itemize}

\subsection{Cross-Modal Generalization Studies (n=5)}

[PLACEHOLDER: Analysis of inter-modality transfer]
\begin{itemize}
    \item Source and target modalities
    \item Reason for multi-modal approach (improved robustness, accessibility, etc.)
    \item Transfer performance
    \item Practical implications
\end{itemize}

